\documentclass[11pt]{article}
\usepackage[margin=1in]{geometry}
\usepackage{graphicx}
\usepackage{booktabs}
\usepackage{array}
\usepackage{longtable}
\usepackage{hyperref}
\usepackage{amsmath}
\usepackage{amssymb}
\usepackage{siunitx}
\hypersetup{colorlinks=true,linkcolor=blue,citecolor=blue,urlcolor=blue}
\title{Conformal \(U(1)_D\) Dark-Sector First-Order Phase Transition and the NANOGrav Signal}
\author{DarkAgents}
\date{8 June 2026}
\begin{document}
\maketitle

\begin{abstract}
We report the DarkAgents analysis of the prompt: ``Can a conformal U(1) BSM extension with a dark scalar, a dark gauge bosons and a dark fermion explain the NANOGrav signal with a cosmological FOPT?''  The conditional answer is yes at the level computed here: the validated conformal Abelian dark sector has viable first-order phase-transition points, a benchmark scan contains points lying inside all 14 PTA violin windows, and a PTArcade run on NANOGrav 15-year data finds a posterior region near \(v_D\simeq 1\) GeV and \(g_D\simeq 0.29\).  This is not yet a full exclusion-safe cosmological model.  The result depends on an approximate semianalytic FOPT backend, a DBF gravitational-wave template used outside strict validity for strongly supercooled points, exact-zero portal assumptions, and missing BBN/CMB, relic-density, thermal-history, RG-validity and completion calculations.
\end{abstract}

\section{Model}
The validated model is the Standard Model plus a dark \(U(1)_D\).  All new fields are SM singlets.  The dark scalar \(\Phi\) has \(U(1)_D\) charge \(2\), the dark gauge boson is \(X_\mu\), and the fermion sector contains two left-handed Weyl fields \(\psi_+\) and \(\psi_-\) with charges \(+1\) and \(-1\).  This charge assignment permits Majorana-type Yukawa terms after \(\Phi\) gets a Coleman-Weinberg vev and is anomaly-safe for the Abelian dark gauge symmetry.

The dark-sector Lagrangian terms used by the validated handoff are
\begin{align}
\mathcal{L}_D \supset
-\frac14 X_{\mu\nu}X^{\mu\nu}
 + |D_\mu\Phi|^2
 + i\psi_+^\dagger\bar\sigma^\mu D_\mu\psi_+
 + i\psi_-^\dagger\bar\sigma^\mu D_\mu\psi_-
 -\frac12 y_D \Phi^\dagger \psi_+\psi_+
 -\frac12 y_D \Phi \psi_-\psi_- + \mathrm{h.c.}
 - \lambda_D(\Phi^\dagger\Phi)^2 .
\end{align}
The symmetry-allowed Higgs portal \(\lambda_{HP}(H^\dagger H)(\Phi^\dagger\Phi)\) and hypercharge kinetic mixing \(\epsilon B_{\mu\nu}X^{\mu\nu}/2\) are fixed to zero in the backend.  The independent scan parameters are \(v_D\), \(g_D\), and \(y_D\).  The dependent quantities are
\begin{equation}
g_b=2g_D,\qquad g_f=\frac{y_D}{\sqrt{2}},\qquad
\beta_\lambda=\frac{3g_b^4-4g_f^4}{8\pi^2},
\end{equation}
with Coleman-Weinberg scalar mass \(m_\chi^2=\beta_\lambda v_D^2\).  The field-dependent masses used for the backend are
\begin{equation}
m_X^2(\chi)=4g_D^2\chi^2,\qquad
m_\psi^2(\chi)=\frac{y_D^2}{2}\chi^2,
\end{equation}
where the fermion entry represents two degenerate Majorana Weyl mass eigenstates with total fermionic degree count \(4\).  The explicit thermal-mass formulae in the handoff are documentation placeholders, not direct semianalytic backend inputs.

\section{Pipeline Status}
\begin{center}
\begin{tabular}{p{0.25\linewidth}p{0.16\linewidth}p{0.49\linewidth}}
\toprule
Stage & Status & Interpretation \\
\midrule
Proposal & warning & Usable conformal \(U(1)_D\) model proposal; no numerical FOPT/PTA claim. \\
Preliminary literature & ok & Relevant conformal dark-sector PTA literature and benchmark seeds identified. \\
Critique/model validation & warning & No fatal model red flags; exact-zero portals, \(\beta_\lambda>0\), and backend caveats remain. \\
FOPT & ok & Semianalytic scan executed. \\
PTA/PTArcade & warning & Usable benchmark and PTArcade outputs with DBF-template and MAP-boundary warnings. \\
Constraints & warning & No non-PTA published constraint applied as an exclusion; several require new observables. \\
Prior audit & warning & Highest-severity unresolved issues are thermal history, portal assumptions, supercooling completion, and physical posterior masking. \\
\bottomrule
\end{tabular}
\end{center}

\section{Literature Context}
The preliminary literature search found that the broad model class is not new: conformal or nearly conformal dark sectors have already been studied as PTA/NANOGrav FOPT explanations.  The exact backend limit used here, with \(q_\Phi=2\), an anomaly-safe Weyl pair, \(\epsilon=0\), and \(\lambda_{HP}=0\), is a variant of that class rather than a fully distinct unexplored mechanism.

The closest verified paper is Balan, Bringmann, Kahlhoefer, Matuszak and Tasillo, ``Sub-GeV dark matter and nano-Hertz gravitational waves from a classically conformal dark sector,'' arXiv:2502.19478.  It contains a conformal hidden \(U(1)'\), dark photon, dark Higgs, stable fermion, Yukawa coupling, and PTA/NANOGrav-relevant FOPT.  Differences include charge normalization and nonzero portal/kinetic-mixing assumptions.  Borah, Dasgupta and Kang, arXiv:2105.01007, studied a dark \(U(1)_D\) FOPT in light of NANOGrav 12.5-year data, but with negligible dark Yukawa in the FOPT analysis.  Goncalves, Marfatia, Morais and Pasechnik, arXiv:2501.11619, studied supercooled conformal dark \(U(1)'\) transitions with PTArcade likelihoods.  Fujikura, Girmohanta, Nakai and Suzuki, arXiv:2306.17086, is relevant to conformal dark PTA signals but uses a different confining mechanism.  Hashino, Kanemura, Takahashi and Tanaka provide Abelian conformal phase-transition context through an INSPIRE-verified record, but were not used as a direct target-model benchmark.

\section{FOPT Results}
The FOPT agent used the local semianalytic conformal single-field backend.  It scanned 791 total points: 416 were viable, 307 failed numerically, and 68 were excluded or physically failed.  Among viable points, 187 had estimated peak-frequency scale in the PTA band.  The backend caveats are important: it uses a small-field/high-temperature polynomial approximation, approximate percolation treatment, and no RG running, portal coupling, kinetic mixing, or full daisy-resummed finite-temperature potential.

\begin{center}
\begin{tabular}{lrrrrrr}
\toprule
Point & \(v_D\) [GeV] & \(g_D\) & \(y_D\) & \(\alpha\) & \(\beta/H\) & \(T_{\rm reh}\) [GeV] \\
\midrule
extended\_low\_v0.002\_g0.5\_y0.5 & 0.002 & 0.50 & 0.50 & 0.804 & 253.7 & \(3.53\times10^{-4}\) \\
extended\_broad\_v0.08367\_g0.25\_y0.25 & 0.0837 & 0.25 & 0.25 & \(2.04\times10^{16}\) & 13.0 & \(6.68\times10^{-3}\) \\
extended\_low\_v0.01108\_g0.38\_y0.5 & 0.0111 & 0.38 & 0.50 & 83.7 & 80.0 & \(1.19\times10^{-3}\) \\
\bottomrule
\end{tabular}
\end{center}
The quantity \(f_{\rm peak,est}\) in the handoff is an order-of-magnitude redshifted scale guide used to seed PTA analysis, not a true spectral peak.

\section{PTA and PTArcade Results}
The PTA agent used a DBF spectrum template.  In the benchmark comparison, 19 viable points were found by the PTA scoring procedure and the best points scored \(14/14\), meaning their spectra fell inside all 14 tabulated PTA violin windows.  The top point has \(v_D=0.7653\) GeV, \(g_D=0.35\), \(y_D=0.9\), \(\alpha=4.22\times10^{11}\), \(\beta/H=17.52\), \(T_\star=0.0469\) GeV, peak frequency \(1.14\times10^{-8}\) Hz, and peak \(h^2\Omega_{\rm GW}=4.80\times10^{-9}\).  This is strong evidence of spectral compatibility with the PTA windows, but the strongest points have extremely large \(\alpha\), so the DBF template validity warning is central.

The PTArcade run used NANOGrav 15-year data in ceffyl mode.  The priors were log-uniform \(v_D\in[10^{-3.5},10]\) GeV, uniform \(g_D\in[0.15,1.0]\), and uniform \(y_D\in[0,0.95]\).  The stored chains were finite; 10001 rows were stored after a completed \(N_{\rm samples}=100000\) run.  The posterior estimates are
\begin{align}
v_D &= 0.960 \ {\rm GeV}\quad(\sigma_{\log_{10}v_D}=0.182),\\
g_D &= 0.286 \pm 0.026,\qquad
y_D = 0.405 \pm 0.255 .
\end{align}
The 68\% intervals are \(v_D=0.652\)--\(1.395\) GeV, \(g_D=0.260\)--\(0.315\), and \(y_D=0.119\)--\(0.705\).  The MAP point is \(v_D=1.176\) GeV, \(g_D=0.262\), \(y_D=0\).  Since \(y_D\) is at the physical lower boundary, the MAP must not be interpreted as an interior preference for a vanishing Yukawa.

\begin{figure}[ht]
\centering
\includegraphics[width=0.78\linewidth]{pta_benchmark_spectrum.pdf}
\caption{Benchmark DBF spectrum compared with the PTA violin windows.}
\end{figure}
\begin{figure}[ht]
\centering
\includegraphics[width=0.78\linewidth]{pta_ptarcade_spectrum.pdf}
\caption{PTArcade MAP spectrum generated from the final chain summary.}
\end{figure}
\begin{figure}[ht]
\centering
\includegraphics[width=0.78\linewidth]{pta_posteriors.pdf}
\caption{PTArcade posterior distributions for the model parameters.}
\end{figure}

\section{Constraints}
The constraint handoff considered 12 constraints.  None was directly testable as a published exclusion with the present observables.  Two were approximately testable, two require a recast, four require new backend calculations, and four were not applicable in the exact-zero-portal backend limit.

\begin{longtable}{p{0.16\linewidth}p{0.23\linewidth}p{0.18\linewidth}p{0.31\linewidth}}
\toprule
ID & Constraint & Status & Reason \\
\midrule
C01 & BBN expansion/dark radiation & requires new backend & Needs dark-sector temperature ratio, decoupling, abundances, and lifetimes. \\
C03 & CMB \(\Delta N_{\rm eff}\) & requires new backend & Radiation density at BBN/recombination was not computed. \\
C04 & CMB energy injection & recast needed & Requires abundance, annihilation/decay channels, and SM-visible final states. \\
C05 & Thermalization/decoupling & requires new backend & Zero portals leave the dark-sector temperature history unspecified. \\
C06--C08 & Dark photon and scalar laboratory limits & not applicable & Published limits require nonzero kinetic mixing or Higgs portal; backend fixes both to zero. \\
C09 & Dark fermion relic abundance & requires new backend & Stability and Boltzmann abundance are not established. \\
C10 & Self-interactions/structure formation & recast needed & Requires DM identity, relic fraction, and velocity-dependent cross sections. \\
C11 & Supercooling completion & approximate & Large-\(\alpha\) points need stronger completion and wall-velocity diagnostics. \\
C12 & Perturbativity/CW validity & approximate & Couplings are perturbative, but rectangular 68\% posterior corners include \(\beta_\lambda<0\); RG running was not computed. \\
\bottomrule
\end{longtable}
The preferred 68\% PTArcade region maps to \(m_X=0.340\)--\(0.878\) GeV and \(m_\psi=0.055\)--\(0.695\) GeV under \(m_X=2g_Dv_D\) and \(m_\psi=y_Dv_D/\sqrt2\).  For \(\beta_\lambda>0\) corners, the estimated scalar mass range is \(m_\chi\simeq0.034\)--\(0.108\) GeV.  Because these ranges use independent rectangular intervals rather than a masked posterior covariance, they should be treated as approximate.

\section{Prior and Assumption Audit}
The highest-severity uncontrolled issues are:
\begin{center}
\begin{tabular}{p{0.08\linewidth}p{0.13\linewidth}p{0.58\linewidth}p{0.10\linewidth}}
\toprule
ID & Category & Issue & Severity \\
\midrule
A01 & portal & Exact \(\epsilon=0\) removes visible dark-photon production, decay, and thermalization channels without a protective argument. & 8 \\
A02 & portal & Exact \(\lambda_{HP}=0\) removes scalar mixing and SM thermal contact without a protective argument. & 8 \\
A05 & physical mask & The rectangular PTArcade 68\% region contains \(\beta_\lambda<0\) corners. & 8 \\
A07 & percolation & Strongly supercooled points need robust completion, reheating, and wall-velocity diagnostics. & 8 \\
A09 & thermal history & No BBN/CMB/\(\Delta N_{\rm eff}\) thermal-history calculation is present. & 8 \\
A08 & GW template & The DBF template is the closest supported template, not a strict template for all very strong transitions. & 7 \\
A10 & dark matter & Fermion stability and relic abundance are unresolved. & 7 \\
A12 & RG validity & No RG, Landau-pole, or loop-stability audit was performed. & 7 \\
\bottomrule
\end{tabular}
\end{center}
The highest-priority followups are a \(\beta_\lambda>0\)-masked posterior summary, a dark-sector thermal-history module, a full completion/percolation cross-check of the preferred and \(14/14\) points, justification or relaxation of exact-zero portals, and a dark fermion stability/relic-density calculation.

\section{Conclusion}
Within the implemented workflow, the model can conditionally explain the NANOGrav signal through a cosmological first-order phase transition.  This statement is supported by a successful semianalytic FOPT scan, PTA benchmark spectra that include \(14/14\) violin-window matches, and a finite PTArcade chain whose posterior lies near \(v_D\sim1\) GeV and \(g_D\sim0.3\).  The result should be read as positive evidence for spectral compatibility, not as a completed cosmological viability proof.

The unresolved caveats are not cosmetic.  Exact-zero portals control whether visible-sector laboratory, BBN and CMB constraints apply; the dark-sector temperature history is absent; the strongest PTA-compatible points are supercooled and require completion diagnostics beyond the current approximation; the DBF template is only a closest supported template; the posterior must be masked to enforce \(\beta_\lambda>0\); and the dark fermion has not been promoted to a fully calculated relic.  No non-PTA constraint currently excludes the reported preferred region, but this is a non-exclusion due to missing observables and backend scope, not a proof that the region is safe.

\section*{References}
\begin{enumerate}
\item Debasish Borah, Arnab Dasgupta, Sin Kyu Kang, ``Gravitational waves from a dark \(U(1)_D\) phase transition in the light of NANOGrav 12.5 yr data,'' arXiv:2105.01007, Phys. Rev. D 104, 063501 (2021).
\item Joao Goncalves, Danny Marfatia, Antonio P. Morais, Roman Pasechnik, ``Supercooled phase transitions in conformal dark sectors explain NANOGrav data,'' arXiv:2501.11619, Phys. Lett. B 869, 139829 (2025).
\item Sowmiya Balan, Torsten Bringmann, Felix Kahlhoefer, Jonas Matuszak, Carlo Tasillo, ``Sub-GeV dark matter and nano-Hertz gravitational waves from a classically conformal dark sector,'' arXiv:2502.19478, JCAP 08 (2025) 062.
\item Kohei Fujikura, Sudhakantha Girmohanta, Yuichiro Nakai, Motoo Suzuki, ``NANOGrav Signal from a Dark Conformal Phase Transition,'' arXiv:2306.17086.
\item Katsuya Hashino, Shinya Kanemura, Yuta Takahashi, Masatoshi Tanaka, ``Phase transition and vacuum stability in the classically conformal B-L model,'' INSPIRE record 1705442.
\item Marco Hufnagel, Kai Schmidt-Hoberg, Sebastian Wild, ``BBN constraints on MeV-scale dark sectors. Part I. Sterile decays,'' arXiv:1712.03972.
\item Marco Hufnagel, Kai Schmidt-Hoberg, Sebastian Wild, ``BBN constraints on MeV-scale dark sectors. Part II. Electromagnetic decays,'' arXiv:1808.09324.
\item Eleonora Di Valentino, Alessandro Melchiorri, Joseph Silk, ``Cosmological constraints in extended parameter space from the Planck 2018 Legacy release,'' arXiv:1908.01391.
\item Torsten Bringmann, Felix Kahlhoefer, Kai Schmidt-Hoberg, Parampreet Walia, ``Strong constraints on self-interacting dark matter with light mediators,'' arXiv:1612.00845.
\item M. Battaglieri et al., ``US Cosmic Visions: New Ideas in Dark Matter 2017: Community Report,'' arXiv:1707.04591.
\item Marco Fabbrichesi, Emidio Gabrielli, Gaia Lanfranchi, ``The Dark Photon,'' arXiv:2005.01515.
\item Thomas Hambye, Laurent Vanderheyden, ``Minimal self-interacting dark matter models with light mediator,'' arXiv:1912.11708.
\item Ryusuke Jinno, Thomas Konstandin, Hyeonseok Rubira, Jorinde van de Vis, ``Supercool subtleties of cosmological phase transitions,'' arXiv:2212.07559.
\item Alexander Farzinnia, Hong-Jian He, Jing Ren, ``Classically conformal B-L extended Standard Model and Higgs vacuum stability,'' arXiv:1205.2893.  This citation is marked uncertain in the constraint handoff metadata and is not used as an exclusion claim here.
\end{enumerate}

\end{document}
